% !TEX root = ./TFGTeXiS.tex

% ----------------------------------------------------------------------
%
%                            TFMTesis.tex
%
%----------------------------------------------------------------------
%
% Este fichero contiene el "documento maestro" del documento. Lo único
% que hace es configurar el entorno LaTeX e incluir los ficheros .tex
% que contienen cada sección.
%
%----------------------------------------------------------------------
%
% Los ficheros necesarios para este documento son:
%
%       TeXiS/* : ficheros de la plantilla TeXiS.
%       Cascaras/* : ficheros con las partes del documento que no
%          son capítulos ni apéndices (portada, agradecimientos, etc.)
%       Capitulos/*.tex : capítulos de la tesis
%       Apendices/*.tex: apéndices de la tesis
%       constantes.tex: constantes LaTeX
%       config.tex : configuración de la "compilación" del documento
%       guionado.tex : palabras con guiones
%
% Para la bibliografía, además, se necesitan:
%
%       *.bib : ficheros con la información de las referencias
%
% ---------------------------------------------------------------------

\documentclass[12pt,a4paper,twoside]{book}

%
% Definimos  el   comando  \compilaCapitulo,  que   luego  se  utiliza
% (opcionalmente) en config.tex. Quedaría  mejor si también se definiera
% en  ese fichero,  pero por  el modo  en el  que funciona  eso  no es
% posible. Puedes consultar la documentación de ese fichero para tener
% más  información. Definimos también  \compilaApendice, que  tiene el
% mismo  cometido, pero  que se  utiliza para  compilar  únicamente un
% apéndice.
%
%
% Si  queremos   compilar  solo   una  parte  del   documento  podemos
% especificar mediante  \includeonly{...} qué ficheros  son los únicos
% que queremos  que se incluyan.  Esto  es útil por  ejemplo para sólo
% compilar un capítulo.
%
% El problema es que todos aquellos  ficheros que NO estén en la lista
% NO   se  incluirán...  y   eso  también   afecta  a   ficheros  de
% la plantilla...
%
% Total,  que definimos  una constante  con los  ficheros  que siempre
% vamos a querer compilar  (aquellos relacionados con configuración) y
% luego definimos \compilaCapitulo.
\newcommand{\ficherosBasicosTeXiS}{%
TeXiS/TeXiS_pream,TeXiS/TeXiS_cab,TeXiS/TeXiS_bib,TeXiS/TeXiS_cover%
}
\newcommand{\ficherosBasicosTexto}{%
constantes,guionado,Cascaras/bibliografia,config%
}
\newcommand{\compilaCapitulo}[1]{%
\includeonly{\ficherosBasicosTeXiS,\ficherosBasicosTexto,Capitulos/#1}%
}

\newcommand{\compilaApendice}[1]{%
\includeonly{\ficherosBasicosTeXiS,\ficherosBasicosTexto,Apendices/#1}%
}

%- - - - - - - - - - - - - - - - - - - - - - - - - - - - - - - - - - -
%            Preámbulo del documento. Configuraciones varias
%- - - - - - - - - - - - - - - - - - - - - - - - - - - - - - - - - - -

% Define  el  tipo  de  compilación que  estamos  haciendo.   Contiene
% definiciones  de  constantes que  cambian  el  comportamiento de  la
% compilación. Debe incluirse antes del paquete TeXiS/TeXiS.sty
%---------------------------------------------------------------------
%
%                          config.tex
%
%---------------------------------------------------------------------
%
% Contiene la  definición de constantes  que determinan el modo  en el
% que se compilará el documento.
%
%---------------------------------------------------------------------
%
% En concreto, podemos  indicar si queremos "modo release",  en el que
% no  aparecerán  los  comentarios  (creados  mediante  \com{Texto}  o
% \comp{Texto}) ni los "por  hacer" (creados mediante \todo{Texto}), y
% sí aparecerán los índices. El modo "debug" (o mejor dicho en modo no
% "release" muestra los índices  (construirlos lleva tiempo y son poco
% útiles  salvo  para   la  versión  final),  pero  sí   el  resto  de
% anotaciones.
%
% Si se compila con LaTeX (no  con pdflatex) en modo Debug, también se
% muestran en una esquina de cada página las entradas (en el índice de
% palabras) que referencian  a dicha página (consulta TeXiS_pream.tex,
% en la parte referente a show).
%
% El soporte para  el índice de palabras en  TeXiS es embrionario, por
% lo  que no  asumas que  esto funcionará  correctamente.  Consulta la
% documentación al respecto en TeXiS_pream.tex.
%
%
% También  aquí configuramos  si queremos  o  no que  se incluyan  los
% acrónimos  en el  documento final  en la  versión release.  Para eso
% define (o no) la constante \acronimosEnRelease.
%
% Utilizando \compilaCapitulo{nombre}  podemos también especificar qué
% capítulo(s) queremos que se compilen. Si no se pone nada, se compila
% el documento  completo.  Si se pone, por  ejemplo, 01Introduccion se
% compilará únicamente el fichero Capitulos/01Introduccion.tex
%
% Para compilar varios  capítulos, se separan sus nombres  con comas y
% no se ponen espacios de separación.
%
% En realidad  la macro \compilaCapitulo  está definida en  el fichero
% principal tesis.tex.
%
%---------------------------------------------------------------------


% Comentar la línea si no se compila en modo release.
% TeXiS hará el resto.
% ¡¡¡Si cambias esto, haz un make clean antes de recompilar!!!
\def\release{1}


% Descomentar la linea si se quieren incluir los
% acrónimos en modo release (en modo debug
% no se incluirán nunca).
% ¡¡¡Si cambias esto, haz un make clean antes de recompilar!!!
%\def\acronimosEnRelease{1}


% Descomentar la línea para establecer el capítulo que queremos
% compilar

% \compilaCapitulo{01Introduccion}
% \compilaCapitulo{02EstructuraYGeneracion}
% \compilaCapitulo{03Edicion}
% \compilaCapitulo{04Imagenes}
% \compilaCapitulo{05Bibliografia}
% \compilaCapitulo{06Makefile}

% \compilaApendice{01AsiSeHizo}

% Variable local para emacs, para  que encuentre el fichero maestro de
% compilación y funcionen mejor algunas teclas rápidas de AucTeX
%%%
%%% Local Variables:
%%% mode: latex
%%% TeX-master: "./Tesis.tex"
%%% End:


% Paquete de la plantilla
\usepackage{TeXiS/TeXiS}

% Incluimos el fichero con comandos de constantes
%---------------------------------------------------------------------
%
%                          constantes.tex
%
%---------------------------------------------------------------------
%
% Fichero que  declara nuevos comandos LaTeX  sencillos realizados por
% comodidad en la escritura de determinadas palabras
%
%---------------------------------------------------------------------

%%%%%%%%%%%%%%%%%%%%%%%%%%%%%%%%%%%%%%%%%%%%%%%%%%%%%%%%%%%%%%%%%%%%%%
% Comando: 
%
%       \titulo
%
% Resultado: 
%
% Escribe el título del documento.
%%%%%%%%%%%%%%%%%%%%%%%%%%%%%%%%%%%%%%%%%%%%%%%%%%%%%%%%%%%%%%%%%%%%%%
\def\titulo{Herramienta de composición de música para videojuegos basada en IA}

%%%%%%%%%%%%%%%%%%%%%%%%%%%%%%%%%%%%%%%%%%%%%%%%%%%%%%%%%%%%%%%%%%%%%%
% Comando: 
%
%       \autor
%
% Resultado: 
%
% Escribe el autor del documento.
%%%%%%%%%%%%%%%%%%%%%%%%%%%%%%%%%%%%%%%%%%%%%%%%%%%%%%%%%%%%%%%%%%%%%%
\def\autor{Javier Callejo Herrero, V\'ictor Manuel Estremera Herranz, Miguel Gonz\'alez P\'erez y Rodrigo S\'anchez Torres}

% Variable local para emacs, para  que encuentre el fichero maestro de
% compilación y funcionen mejor algunas teclas rápidas de AucTeX

%%%
%%% Local Variables:
%%% mode: latex
%%% TeX-master: "tesis.tex"
%%% End:
%%%%%%%%%%%%%%%%%%%%%%%%%%%%%%%%%%%%%%%%%%%%%%%%%%%%%%%%%%%%%%%%%%%%%%
% Comando: 
%
%       \appName
%
% Resultado: 
%
% Escribe el nombre de la aplicación
%%%%%%%%%%%%%%%%%%%%%%%%%%%%%%%%%%%%%%%%%%%%%%%%%%%%%%%%%%%%%%%%%%%%%%
\def\appName{Vanguard Music}
%%%%%%%%%%%%%%%%%%%%%%%%%%%%%%%%%%%%%%%%%%%%%%%%%%%%%%%%%%%%%%%%%%%%%%
\def\generationTabName{Generación}
\def\tematicTabName{Musicalización}
\def\advancedTabName{Avanzado}
\def\configTabName{Configuración}

% Sacamos en el log de la compilación el copyright
%\typeout{Copyright Marco Antonio and Pedro Pablo Gomez Martin}

%
% "Metadatos" para el PDF
%
\ifpdf\hypersetup{%
    pdftitle = {\titulo},
    pdfsubject = {Plantilla de Tesis},
    pdfkeywords = {Plantilla, LaTeX, tesis, trabajo de
      investigación, trabajo de Master},
    pdfauthor = {\textcopyright\ \autor},
    pdfcreator = {\LaTeX\ con el paquete \flqq hyperref\frqq},
    pdfproducer = {pdfeTeX-0.\the\pdftexversion\pdftexrevision},
    }
    \pdfinfo{/CreationDate (\today)}
\fi


%- - - - - - - - - - - - - - - - - - - - - - - - - - - - - - - - - - -
%                        Documento
%- - - - - - - - - - - - - - - - - - - - - - - - - - - - - - - - - - -
\begin{document}

% Incluimos el  fichero de definición de guionado  de algunas palabras
% que LaTeX no ha dividido como debería
%----------------------------------------------------------------
%
%                          guionado.tex
%
%----------------------------------------------------------------
%
% Fichero con algunas divisiones de palabras que LaTeX no
% hace correctamente si no se le da alguna ayuda.
%
%----------------------------------------------------------------

\hyphenation{
% a
abs-trac-to
abs-trac-tos
abs-trac-ta
abs-trac-tas
ac-tua-do-res
a-gra-de-ci-mien-tos
ana-li-za-dor
an-te-rio-res
an-te-rior-men-te
apa-rien-cia
a-pro-pia-do
a-pro-pia-dos
a-pro-pia-da
a-pro-pia-das
a-pro-ve-cha-mien-to
a-que-llo
a-que-llos
a-que-lla
a-que-llas
a-sig-na-tu-ra
a-sig-na-tu-ras
a-so-cia-da
a-so-cia-das
a-so-cia-do
a-so-cia-dos
au-to-ma-ti-za-do
% b
batch
bi-blio-gra-fía
bi-blio-grá-fi-cas
bien
bo-rra-dor
boo-l-ean-expr
% c
ca-be-ce-ra
call-me-thod-ins-truc-tion
cas-te-lla-no
cir-cuns-tan-cia
cir-cuns-tan-cias
co-he-ren-te
co-he-ren-tes
co-he-ren-cia
co-li-bri
co-men-ta-rio
co-mer-cia-les
co-no-ci-mien-to
cons-cien-te
con-si-de-ra-ba
con-si-de-ra-mos
con-si-de-rar-se
cons-tan-te
cons-trucción
cons-tru-ye
cons-tru-ir-se
con-tro-le
co-rrec-ta-men-te
co-rres-pon-den
co-rres-pon-dien-te
co-rres-pon-dien-tes
co-ti-dia-na
co-ti-dia-no
crean
cris-ta-li-zan
cu-rri-cu-la
cu-rri-cu-lum
cu-rri-cu-lar
cu-rri-cu-la-res
% d
de-di-ca-do
de-di-ca-dos
de-di-ca-da
de-di-ca-das
de-rro-te-ro
de-rro-te-ros
de-sa-rro-llo
de-sa-rro-llos
de-sa-rro-lla-do
de-sa-rro-lla-dos
de-sa-rro-lla-da
de-sa-rro-lla-das
de-sa-rro-lla-dor
de-sa-rro-llar
des-cri-bi-re-mos
des-crip-ción
des-crip-cio-nes
des-cri-to
des-pués
de-ta-lla-do
de-ta-lla-dos
de-ta-lla-da
de-ta-lla-das
di-a-gra-ma
di-a-gra-mas
di-se-ños
dis-po-ner
dis-po-ni-bi-li-dad
do-cu-men-ta-da
do-cu-men-to
do-cu-men-tos
% e
edi-ta-do
e-du-ca-ti-vo
e-du-ca-ti-vos
e-du-ca-ti-va
e-du-ca-ti-vas
e-la-bo-ra-do
e-la-bo-ra-dos
e-la-bo-ra-da
e-la-bo-ra-das
es-co-llo
es-co-llos
es-tu-dia-do
es-tu-dia-dos
es-tu-dia-da
es-tu-dia-das
es-tu-dian-te
e-va-lua-cio-nes
e-va-lua-do-res
exis-ten-tes
exhaus-ti-va
ex-pe-rien-cia
ex-pe-rien-cias
% f
for-ma-li-za-do
% g
ge-ne-ra-ción
ge-ne-ra-dor
ge-ne-ra-do-res
ge-ne-ran
% h
he-rra-mien-ta
he-rra-mien-tas
% i
i-dio-ma
i-dio-mas
im-pres-cin-di-ble
im-pres-cin-di-bles
in-de-xa-do
in-de-xa-dos
in-de-xa-da
in-de-xa-das
in-di-vi-dual
in-fe-ren-cia
in-fe-ren-cias
in-for-ma-ti-ca
in-gre-dien-te
in-gre-dien-tes
in-me-dia-ta-men-te
ins-ta-la-do
ins-tan-cias
% j
% k
% l
len-gua-je
li-be-ra-to-rio
li-be-ra-to-rios
li-be-ra-to-ria
li-be-ra-to-rias
li-mi-ta-do
li-te-ra-rio
li-te-ra-rios
li-te-ra-ria
li-te-ra-rias
lo-tes
% m
ma-ne-ra
ma-nual
mas-que-ra-de
ma-yor
me-mo-ria
mi-nis-te-rio
mi-nis-te-rios
mo-de-lo
mo-de-los
mo-de-la-do
mo-du-la-ri-dad
mo-vi-mien-to
% n
na-tu-ral
ni-vel
nues-tro
% o
obs-tan-te
o-rien-ta-do
o-rien-ta-dos
o-rien-ta-da
o-rien-ta-das
% p
pa-ra-le-lo
pa-ra-le-la
par-ti-cu-lar
par-ti-cu-lar-men-te
pe-da-gó-gi-ca
pe-da-gó-gi-cas
pe-da-gó-gi-co
pe-da-gó-gi-cos
pe-rio-di-ci-dad
per-so-na-je
plan-te-a-mien-to
plan-te-a-mien-tos
po-si-ción
pre-fe-ren-cia
pre-fe-ren-cias
pres-cin-di-ble
pres-cin-di-bles
pri-me-ra
pro-ble-ma
pro-ble-mas
pró-xi-mo
pu-bli-ca-cio-nes
pu-bli-ca-do
% q
% r
rá-pi-da
rá-pi-do
ra-zo-na-mien-to
ra-zo-na-mien-tos
re-a-li-zan-do
re-fe-ren-cia
re-fe-ren-cias
re-fe-ren-cia-da
re-fe-ren-cian
re-le-van-tes
re-pre-sen-ta-do
re-pre-sen-ta-dos
re-pre-sen-ta-da
re-pre-sen-ta-das
re-pre-sen-tar-lo
re-qui-si-to
re-qui-si-tos
res-pon-der
res-pon-sa-ble
% s
se-pa-ra-do
si-guien-do
si-guien-te
si-guien-tes
si-guie-ron
si-mi-lar
si-mi-la-res
si-tua-ción
% t
tem-pe-ra-ments
te-ner
trans-fe-ren-cia
trans-fe-ren-cias
% u
u-sua-rio
Unreal-Ed
% v
va-lor
va-lo-res
va-rian-te
ver-da-de-ro
ver-da-de-ros
ver-da-de-ra
ver-da-de-ras
ver-da-de-ra-men-te
ve-ri-fi-ca
% w
% x
% y
% z
}
% Variable local para emacs, para que encuentre el fichero
% maestro de compilación
%%%
%%% Local Variables:
%%% mode: latex
%%% TeX-master: "./Tesis.tex"
%%% End:


% Marcamos  el inicio  del  documento para  la  numeración de  páginas
% (usando números romanos para esta primera fase).
\frontmatter
\pagestyle{empty}

%---------------------------------------------------------------------
%
%                          configCover.tex
%
%---------------------------------------------------------------------
%
% cover.tex
% Copyright 2009 Marco Antonio Gomez-Martin, Pedro Pablo Gomez-Martin
%
% This file belongs to the TeXiS manual, a LaTeX template for writting
% Thesis and other documents. The complete last TeXiS package can
% be obtained from http://gaia.fdi.ucm.es/projects/texis/
%
% Although the TeXiS template itself is distributed under the 
% conditions of the LaTeX Project Public License
% (http://www.latex-project.org/lppl.txt), the manual content
% uses the CC-BY-SA license that stays that you are free:
%
%    - to share & to copy, distribute and transmit the work
%    - to remix and to adapt the work
%
% under the following conditions:
%
%    - Attribution: you must attribute the work in the manner
%      specified by the author or licensor (but not in any way that
%      suggests that they endorse you or your use of the work).
%    - Share Alike: if you alter, transform, or build upon this
%      work, you may distribute the resulting work only under the
%      same, similar or a compatible license.
%
% The complete license is available in
% http://creativecommons.org/licenses/by-sa/3.0/legalcode
%
%---------------------------------------------------------------------
%
% Fichero que contiene la configuración de la portada y de la 
% primera hoja del documento.
%
%---------------------------------------------------------------------


% Pueden configurarse todos los elementos del contenido de la portada
% utilizando comandos.

%%%%%%%%%%%%%%%%%%%%%%%%%%%%%%%%%%%%%%%%%%%%%%%%%%%%%%%%%%%%%%%%%%%%%%
% Título del documento:
% \tituloPortada{titulo}
% Nota:
% Si no se define se utiliza el del \titulo. Este comando permite
% cambiar el título de forma que se especifiquen dónde se quieren
% los retornos de carro cuando se utilizan fuentes grandes.
%%%%%%%%%%%%%%%%%%%%%%%%%%%%%%%%%%%%%%%%%%%%%%%%%%%%%%%%%%%%%%%%%%%%%%
\tituloPortada{%
Asistente musical con IA
}


%%%%%%%%%%%%%%%%%%%%%%%%%%%%%%%%%%%%%%%%%%%%%%%%%%%%%%%%%%%%%%%%%%%%%%
% Título del documento en inglés:
% \tituloPortadaEng{titulo}
% Nota:
% Si no se define se utiliza el del \titulo. Este comando permite
% cambiar el título de forma que se especifiquen dónde se quieren
% los retornos de carro cuando se utilizan fuentes grandes.
%%%%%%%%%%%%%%%%%%%%%%%%%%%%%%%%%%%%%%%%%%%%%%%%%%%%%%%%%%%%%%%%%%%%%%
\tituloPortadaEng{%
Title in English \textcolor{red}{(defined in Cascaras$\backslash$cover.tex)}
}

%%%%%%%%%%%%%%%%%%%%%%%%%%%%%%%%%%%%%%%%%%%%%%%%%%%%%%%%%%%%%%%%%%%%%%
% Autor del documento:
% \autorPortada{Nombre}
% Se utiliza en la portada y en el valor por defecto del
% primer subtítulo de la segunda portada.
%%%%%%%%%%%%%%%%%%%%%%%%%%%%%%%%%%%%%%%%%%%%%%%%%%%%%%%%%%%%%%%%%%%%%%
\autorPortada{\textcolor{red}{Javier Callejo Herrero
                            \\Víctor Manuel Estremera Herranz
                            \\Miguel González Pérez
                            \\Rodrigo Sánchez Torres
                            }}

%%%%%%%%%%%%%%%%%%%%%%%%%%%%%%%%%%%%%%%%%%%%%%%%%%%%%%%%%%%%%%%%%%%%%%
% Fecha de publicación:
% \fechaPublicacion{Fecha}
% Puede ser vacío. Aparece en la última línea de ambas portadas
%%%%%%%%%%%%%%%%%%%%%%%%%%%%%%%%%%%%%%%%%%%%%%%%%%%%%%%%%%%%%%%%%%%%%%
% Descomentar para que ponga siempre la fecha actual
\fechaPublicacion{\today}
%\fechaPublicacion{\textcolor{red}{DIA de MES de AÑO}}

%%%%%%%%%%%%%%%%%%%%%%%%%%%%%%%%%%%%%%%%%%%%%%%%%%%%%%%%%%%%%%%%%%%%%%
% Imagen de la portada (y escala)
% \imagenPortada{Fichero}
% \escalaImagenPortada{Numero}
% Si no se especifica, se utiliza la imagen TODO.pdf
%%%%%%%%%%%%%%%%%%%%%%%%%%%%%%%%%%%%%%%%%%%%%%%%%%%%%%%%%%%%%%%%%%%%%%
% imagen en blanco y negro
%\imagenPortada{Imagenes/Vectorial/escudoUCM}
%imagen en color
\imagenPortada{Imagenes/Bitmap/escudoUCMcolor}
\escalaImagenPortada{.2}

%%%%%%%%%%%%%%%%%%%%%%%%%%%%%%%%%%%%%%%%%%%%%%%%%%%%%%%%%%%%%%%%%%%%%%
% Tipo de documento.
% \tipoDocumento{Tipo}
% Para el texto justo debajo del escudo.
% Si no se indica, se utiliza "TESIS DOCTORAL".
%%%%%%%%%%%%%%%%%%%%%%%%%%%%%%%%%%%%%%%%%%%%%%%%%%%%%%%%%%%%%%%%%%%%%%
\tipoDocumento{Trabajo de Fin de Grado}

%%%%%%%%%%%%%%%%%%%%%%%%%%%%%%%%%%%%%%%%%%%%%%%%%%%%%%%%%%%%%%%%%%%%%%
% Institución/departamento asociado al documento.
% \institucion{Nombre}
% Puede tener varias líneas. Se utiliza en las dos portadas.
% Si no se indica aparecerá vacío.
%%%%%%%%%%%%%%%%%%%%%%%%%%%%%%%%%%%%%%%%%%%%%%%%%%%%%%%%%%%%%%%%%%%%%%
\institucion{%
Grado en \textcolor{red}{Desarrollo de Videojuegos}\\[0.2em]
Facultad de Informática\\[0.2em]
Universidad Complutense de Madrid
}

%%%%%%%%%%%%%%%%%%%%%%%%%%%%%%%%%%%%%%%%%%%%%%%%%%%%%%%%%%%%%%%%%%%%%%
% Director del trabajo.
% \directorPortada{Nombre}
% Se utiliza para el valor por defecto del segundo subtítulo, donde
% se indica quién es el director del trabajo.
% Si se fuerza un subtítulo distinto, no hace falta definirlo.
%%%%%%%%%%%%%%%%%%%%%%%%%%%%%%%%%%%%%%%%%%%%%%%%%%%%%%%%%%%%%%%%%%%%%%
\directorPortada{\textcolor{red}{Miguel Gómez-Zamalloa Gil\\Jaime Sánchez Hernández}}


%%%%%%%%%%%%%%%%%%%%%%%%%%%%%%%%%%%%%%%%%%%%%%%%%%%%%%%%%%%%%%%%%%%%%%
% Colaborador en la dirección del trabajo.
% \colaboradorPortada{Nombre}
% Se utiliza para el valor por defecto del segundo subtítulo, donde
% se indica quién es el colaborador en la dirección del trabajo.
% Si se fuerza un subtítulo distinto, no hace falta definirlo.
%%%%%%%%%%%%%%%%%%%%%%%%%%%%%%%%%%%%%%%%%%%%%%%%%%%%%%%%%%%%%%%%%%%%%%
\colaboradorPortada{\textcolor{red}{Colaborador 1\\Colaborador 2}}


%%%%%%%%%%%%%%%%%%%%%%%%%%%%%%%%%%%%%%%%%%%%%%%%%%%%%%%%%%%%%%%%%%%%%%
% Texto del primer subtítulo de la segunda portada.
% \textoPrimerSubtituloPortada{Texto}
% Para configurar el primer "texto libre" de la segunda portada.
% Si no se especifica se indica "Memoria que presenta para optar al
% título de Doctor en Informática" seguido del \autorPortada.
%%%%%%%%%%%%%%%%%%%%%%%%%%%%%%%%%%%%%%%%%%%%%%%%%%%%%%%%%%%%%%%%%%%%%%
\textoPrimerSubtituloPortada{%
\textbf{Trabajo de Fin de Grado en \textcolor{red}{Desarrollo de Videojuegos}}\\ [0.3em]
}

%%%%%%%%%%%%%%%%%%%%%%%%%%%%%%%%%%%%%%%%%%%%%%%%%%%%%%%%%%%%%%%%%%%%%%
% Texto del segundo subtítulo de la segunda portada.
% \textoSegundoSubtituloPortada{Texto}
% Para configurar el segundo "texto libre" de la segunda portada.
% Si no se especifica se indica "Dirigida por el Doctor" seguido
% del \directorPortada.
%%%%%%%%%%%%%%%%%%%%%%%%%%%%%%%%%%%%%%%%%%%%%%%%%%%%%%%%%%%%%%%%%%%%%%
\textoSegundoSubtituloPortada{%
\textbf{Convocatoria: }\textit{\textcolor{red}{Junio} \the\year}%\\[0.2em]
%\textbf{Calificación: }\textit{\textcolor{red}{Nota}}
}

%%%%%%%%%%%%%%%%%%%%%%%%%%%%%%%%%%%%%%%%%%%%%%%%%%%%%%%%%%%%%%%%%%%%%%
% \explicacionDobleCara
% Si se utiliza, se aclara que el documento está preparado para la
% impresión a doble cara.
%%%%%%%%%%%%%%%%%%%%%%%%%%%%%%%%%%%%%%%%%%%%%%%%%%%%%%%%%%%%%%%%%%%%%%
%\explicacionDobleCara

%%%%%%%%%%%%%%%%%%%%%%%%%%%%%%%%%%%%%%%%%%%%%%%%%%%%%%%%%%%%%%%%%%%%%%
% \isbn
% Si se utiliza, aparecerá el ISBN detrás de la segunda portada.
%%%%%%%%%%%%%%%%%%%%%%%%%%%%%%%%%%%%%%%%%%%%%%%%%%%%%%%%%%%%%%%%%%%%%%
%\isbn{978-84-692-7109-4}


%%%%%%%%%%%%%%%%%%%%%%%%%%%%%%%%%%%%%%%%%%%%%%%%%%%%%%%%%%%%%%%%%%%%%%
% \copyrightInfo
% Si se utiliza, aparecerá información de los derechos de copyright
% detrás de la segunda portada.
%%%%%%%%%%%%%%%%%%%%%%%%%%%%%%%%%%%%%%%%%%%%%%%%%%%%%%%%%%%%%%%%%%%%%%
%\copyrightInfo{\autor}


%%
%% Creamos las portadas
%%
\makeCover

% Variable local para emacs, para que encuentre el fichero
% maestro de compilación
%%%
%%% Local Variables:
%%% mode: latex
%%% TeX-master: "../Tesis.tex"
%%% End:

%\include{Cascaras/autorizacion}
% +--------------------------------------------------------------------+
% | Dedication Page (Optional)
% +--------------------------------------------------------------------+

\chapter*{Dedicatoria}

\begin{flushright}
\begin{minipage}[c]{8.5cm}
\flushright{\textit{A Jaime y Miky, por descubrirnos este campo y por su implicación en el trabajo}}
\end{minipage}
\end{flushright}
% +--------------------------------------------------------------------+
% | Acknowledgements Page (Optional)                                   |
% +--------------------------------------------------------------------+

\chapter*{Agradecimientos}

A Guillermo, por el tiempo empleado en hacer estas plantillas. A Adrián, Enrique y Nacho, por sus comentarios para mejorar lo que hicimos. Y a Narciso, a quien no le ha hecho falta el Anillo Único para coordinarnos a todos.












\chapter*{Resumen}

\section*{\tituloPortadaVal}

Este trabajo de fin de grado tiene como propósito crear una apliación de escritorio que permita a los usuarios generar música temática para sus videojuegos de forma sencilla, sin requerir conocimientos musicales previos. El nombre que le hemos dado a la aplicación es \appName{}.

Mediante el uso y estudio de campos como el de la Inteligencia Artificial y el Aprendizaje Automático impulsamos la generación de temas musicales que cumplan las expectativas de la aplicación. \appName{} conecta con \textit{Reaper}, una DAW de terceros que nos permite crear y editar pistas de audio de manera profesional, añadiendo los instrumentos y filtros necesarios para sonorizar la música generada por la IA.


\section*{Palabras clave}
   
\noindent Música para videojuegos, Inteligencia Artificial, Redes Neuronales, Modelos de Markov, Composición automática de música.

   



\begin{otherlanguage}{english}
\chapter*{Abstract}

\section*{\tituloPortadaEngVal}

An abstract in English, half a page long, including the title in English. Below, a list with no more than 10 keywords.


\section*{Keywords}

\noindent Music, Artificial Intelligence, Generation, Neural Networks, Markov Models, Generation, Python




% Si el trabajo se escribe en inglés, comentar esta línea y descomentar
% otra igual que hay justo antes de \end{document}
\end{otherlanguage}

\ifx\generatoc\undefined
\else
%---------------------------------------------------------------------
%
%                          TeXiS_toc.tex
%
%---------------------------------------------------------------------
%
% TeXiS_toc.tex
% Copyright 2009 Marco Antonio Gomez-Martin, Pedro Pablo Gomez-Martin
%
% This file belongs to TeXiS, a LaTeX template for writting
% Thesis and other documents. The complete last TeXiS package can
% be obtained from http://gaia.fdi.ucm.es/projects/texis/
%
% This work may be distributed and/or modified under the
% conditions of the LaTeX Project Public License, either version 1.3
% of this license or (at your option) any later version.
% The latest version of this license is in
%   http://www.latex-project.org/lppl.txt
% and version 1.3 or later is part of all distributions of LaTeX
% version 2005/12/01 or later.
%
% This work has the LPPL maintenance status `maintained'.
% 
% The Current Maintainers of this work are Marco Antonio Gomez-Martin
% and Pedro Pablo Gomez-Martin
%
%---------------------------------------------------------------------
%
% Contiene  los  comandos  para  generar los  índices  del  documento,
% entendiendo por índices las tablas de contenidos.
%
% Genera  el  índice normal  ("tabla  de  contenidos"),  el índice  de
% figuras y el de tablas. También  crea "marcadores" en el caso de que
% se esté compilando con pdflatex para que aparezcan en el PDF.
%
%---------------------------------------------------------------------


% Primero un poquito de configuración...


% Pedimos que inserte todos los epígrafes hasta el nivel \subsection en
% la tabla de contenidos.
\setcounter{tocdepth}{2} 

% Le  pedimos  que nos  numere  todos  los  epígrafes hasta  el  nivel
% \subsubsection en el cuerpo del documento.
\setcounter{secnumdepth}{3} 


% Creamos los diferentes índices.

% Lo primero un  poco de trabajo en los marcadores  del PDF. No quiero
% que  salga una  entrada  por cada  índice  a nivel  0...  si no  que
% aparezca un marcador "Índices", que  tenga dentro los otros tipos de
% índices.  Total, que creamos el marcador "Índices".
% Antes de  la creación  de los índices,  se añaden los  marcadores de
% nivel 1.

\ifpdf
   \pdfbookmark{Índices}{indices}
\fi

% Tabla de contenidos.
%
% La  inclusión  de '\tableofcontents'  significa  que  en la  primera
% pasada  de  LaTeX  se  crea   un  fichero  con  extensión  .toc  con
% información sobre la tabla de contenidos (es conceptualmente similar
% al  .bbl de  BibTeX, creo).  En la  segunda ejecución  de  LaTeX ese
% documento se utiliza para  generar la verdadera página de contenidos
% usando la  información sobre los  capítulos y demás guardadas  en el
% .toc
\ifpdf
   \pdfbookmark[1]{Tabla de Contenidos}{tabla de contenidos}
\fi

\cabeceraEspecial{\'Indice}

\tableofcontents

\newpage 

% Índice de figuras
%
% La idea es semejante que para  el .toc del índice, pero ahora se usa
% extensión .lof (List Of Figures) con la información de las figuras.

\ifpdf
   \pdfbookmark[1]{Índice de figuras}{indice de figuras}
\fi

\cabeceraEspecial{\'Indice de figuras}

\listoffigures

\newpage

% Índice de tablas
% Como antes, pero ahora .lot (List Of Tables)

\ifpdf
   \pdfbookmark[1]{Índice de tablas}{indice de tablas}
\fi

\cabeceraEspecial{\'Indice de tablas}

\listoftables

\newpage

% Variable local para emacs, para  que encuentre el fichero maestro de
% compilación y funcionen mejor algunas teclas rápidas de AucTeX

%%%
%%% Local Variables:
%%% mode: latex
%%% TeX-master: "../Tesis.tex"
%%% End:

\fi

% Marcamos el  comienzo de  los capítulos (para  la numeración  de las
% páginas) y ponemos la cabecera normal
\mainmatter

\pagestyle{fancy}
\restauraCabecera


\chapter{Introducción}
\label{cap:introduccion}

\chapterquote{Frase célebre dicha por alguien inteligente}{Autor}

Según la normativa  para Trabajos de Fin de Grado\footnote{\url{https://informatica.ucm.es/file/normativatfg-2021-2022?ver} (ver versión actualizada para cada curso académico)}, la memoria incluirá una portada normalizada con la siguiente información: título en castellano, título en inglés, autores, profesor director, codirector si es el caso, curso
académico e identificación de la asignatura (Trabajo de fin de grado del Grado en -
nombre del grado correspondiente-, Facultad de Informática, Universidad Complutense de Madrid). Los datos referentes al título y director (y codirector en su caso) deben corresponder a los publicados en la página web de TFG.

La memoria debe incluir la descripción detallada de la propuesta hardware/software realizada y ha de contener:
\begin{enumerate}
\renewcommand{\theenumi}{\alph{enumi}}
	\item un índice,
	\item un resumen y una lista de no más de 10 palabras clave para su búsqueda bibliográfica, ambos en castellano e inglés,
	\item una introducción con los antecedentes, objetivos y plan de trabajo,
	\item resultados y discusión crítica y razonada de los mismos, con sus conclusiones,
	\item bibliografía.
\end{enumerate}
Para facilitar la escritura de la memoria siguiendo esta estructura, el estudiante podrá usar las plantillas en LaTeX o Word preparadas al efecto y publicadas en la página web de TFG.

La memoria constará de un mínimo de 25 páginas para los proyectos realizados por un único estudiante, y de al menos 5 páginas más por cada integrante adicional del grupo. En este número de páginas solo se tiene en cuenta el contenido correspondiente a los apartados c y d del punto anterior.

La memoria puede estar escrita en castellano o inglés, pero en el primer caso la introducción y las conclusiones deben aparecer también en inglés. Las memorias de los TFG matriculados en el grupo I deberán estar escritas íntegramente en inglés, excepto por lo especificado en los puntos 1 y 2 anteriores (título, resumen y lista de palabras clave).

En caso de trabajos no unipersonales, cada participante indicará en la memoria su contribución al proyecto con una extensión de al menos dos páginas por cada uno de los participantes.

Todo el material no original, ya sea texto o figuras, deberá ser convenientemente citado y referenciado. En el caso de material complementario se deben respetar las licencias y \emph{copyrights} asociados al software y hardware que se emplee. En caso contrario no se autorizará la defensa, sin menoscabo de otras acciones que correspondan.


\section{Motivación}
Introducción al tema del TFG.


\section{Objetivos}
Descripción de los objetivos del trabajo.


\section{Plan de trabajo}
Aquí se describe el plan de trabajo a seguir para la consecución de los objetivos descritos en el apartado anterior.



\section{Explicaciones adicionales sobre el uso de esta plantilla}
Si quieres cambiar el \textbf{estilo del título} de los capítulos del documento, edita el fichero \verb|TeXiS\TeXiS_pream.tex| y comenta la línea \verb|\usepackage[Lenny]{fncychap}| para dejar el estilo básico de \LaTeX.

Si no te gusta que no haya \textbf{espacios entre párrafos} y quieres dejar un pequeño espacio en blanco, no metas saltos de línea (\verb|\\|) al final de los párrafos. En su lugar, busca el comando  \verb|\setlength{\parskip}{0.2ex}| en \verb|TeXiS\TeXiS_pream.tex| y aumenta el valor de $0.2ex$ a, por ejemplo, $1ex$.

TFGTeXiS se ha elaborado a partir de la plantilla de TeXiS\footnote{\url{http://gaia.fdi.ucm.es/research/texis/}}, creada por Marco Antonio y Pedro Pablo Gómez Martín para escribir su tesis doctoral. Para explicaciones más extensas y detalladas sobre cómo usar esta plantilla, recomendamos la lectura del documento \texttt{TeXiS-Manual-1.0.pdf} que acompaña a esta plantilla.

El siguiente texto se genera con el comando \verb|\lipsum[2-20]| que viene a continuación en el fichero .tex. El único propósito es mostrar el aspecto de las páginas usando esta plantilla. Quita este comando y, si quieres, comenta o elimina el paquete \textit{lipsum} al final de \verb|TeXiS\TeXiS_pream.tex|

\subsection{Texto de prueba}


\lipsum[2-20]
\chapter{Estado de la Cuestión}
\label{cap:estadoDeLaCuestion}

En el estado de la cuestión es donde aparecen gran parte de las referencias bibliográficas del trabajo. Una de las formas más cómodas de gestionar la bibliografía en {\LaTeX} es utilizando \textbf{bibtex}. Las entradas bibliográficas deben estar en un fichero con extensión \textit{.bib} (con esta plantilla se proporciona el fichero biblio.bib, donde están las entradas referenciadas más abajo). Cada entrada bibliográfica tiene una clave que permite referenciarla desde cualquier parte del texto con los siguiente comandos:

\begin{itemize}
\item Referencia bibliografica con cite: \cite{ldesc2e}
\item Referencia bibliográfica con citep: \citep{notsoshort}
\item Referencia bibliográfica con citet: \citet{latexAPrimer}
\item Referencia inventada: \citep{}
\end{itemize}

Es posible citar más de una fuente, como por ejemplo \citep{latexCompanion,LaTeXLamport,texKnuth}

Después, \LaTeX se ocupa de rellenar la sección de bibliografía con las entradas \textbf{que hayan sido citadas} (es decir, no con todas las entradas que hay en el .bib, sino sólo con aquellas que se hayan citado en alguna parte del texto).

Bibtex es un programa separado de latex, pdflatex o cualquier otra cosa que se use para compilar los .tex, de manera que para que se rellene correctamente la sección de bibliografía es necesario compilar primero el trabajo (a veces es necesario compilarlo dos veces), compilar después con bibtex, y volver a compilar otra vez el trabajo (de nuevo, puede ser necesario compilarlo dos veces). 

\chapter{Descripción del Trabajo}
\label{cap:descripcionTrabajo}

Aquí comienza la descripción del trabajo realizado. Se deben incluir tantos capítulos como sea necesario para describir de la manera más completa posible el trabajo que se ha llevado a cabo. Como muestra la figura \ref{fig:sampleImage}, está todo por hacer.

\begin{figure}[h]
	\centering
	\includegraphics[width = 0.5\textwidth]{Imagenes/Vectorial/Todo.pdf}
	\caption{Ejemplo de imagen}
	\label{fig:sampleImage}
\end{figure}

Si te sirve de utilidad,  puedes incluir tablas para mostrar resultados, tal como se ve en la tabla \ref{tab:sampleTable}.

\begin{table}
	\centering
	\begin{tabular}{c|c|c}
		\textbf{Col 1} & \textbf{Col 2} & \textbf{Col 3} \\
		\hline\hline
		3 & 3.01 & 3.50\\
		6 & 2.12 & 4.40\\
		1 & 3.79 & 5.00\\
		2 & 4.88 & 5.30\\
		4 & 3.50 & 2.90\\
		5 & 7.40 & 4.70\\
		\hline
	\end{tabular}
	\caption{Tabla de ejemplo}
	\label{tab:sampleTable}
\end{table}

%\include{Capitulos/Capitulo4}
%\include{Capitulos/Capitulo5}
\chapter{Conclusiones y Trabajo Futuro}
\label{cap:conclusiones}

Conclusiones del trabajo y líneas de trabajo futuro.

Antes de la entrega de actas de cada convocatoria, en el plazo que se indica en el calendario de los trabajos de fin de grado, el estudiante entregará en el Campus Virtual la versión final de la memoria en PDF.




%%%%%%%%%%%%%%%%%%%%%%%%%%%%%%%%%%%%%%%%%%%%%%%%%%%%%%%%%%%%%%%%%%%%%%%%%%%
% Si el TFG se escribe en inglés, comentar las siguientes líneas 
% porque no es necesario incluir nuevamente las Conclusiones en inglés
\begin{otherlanguage}{english}
\chapter*{Introduction}
\label{cap:introduction}
\addcontentsline{toc}{chapter}{Introduction}

Introduction to the subject area. This chapter contains the translation of Chapter \ref{cap:introduccion}.

\section{Motivación}
Our main purpose for this project is to make music composition, focused primarily on video games, to users with low or non existent musical knowledge.

In order to reach this purpose, we want to rely on the field of Artificial Intelligence, exploring techniques and methodologies to be able to generate something that is interesting.

\section{Objectives}
With our purpose defined, the next steps is to define a series of objectives:
\begin{enumerate}
    \item Generate melodies through the use of different tools, including \textit{Machine Learning} or stochastic algorithms.
    \item Harmonize the melody previously generated in a way that results pleasant to the human ear, either through the use of Artificial Intelligence or own heuristics.
    \item Sound the generated music notes making a property selection of instruments, choosing between pre-defined themes. 
    \item Create an application for \textit{Windows} that allows the user to apply this techniques in a friendly way.
\end{enumerate}

\section{Plan de trabajo}
El tener objetivos claramente separados en cuanto a funcionalidad, nos permite separar el trabajo de manera más eficiente en módulos independientes, como una cadena de producción separada en etapas que se une en la app final. 

\begin{enumerate}
    \item Investigación y estudio del campo de la Inteligencia Artificial aplicada a la generación musical. Durante esta fase, revisaremos la literatura existente, investigaremos los enfoques y algoritmos más relevantes y nos familiarizaremos con las herramientas y recursos disponibles.
    \item Desarrollo del algoritmo de generación de melodías. Utilizaremos técnicas de aprendizaje automático y modelado generativo para crear un sistema capaz de componer melodías originales y variadas.
    \item Implementación de la armonización de las melodías generadas. En esta etapa, diseñaremos y desarrollaremos heurísticas y reglas musicales para armonizar las melodías de manera coherente y agradable al oído.
    \item Selección y configuración de instrumentos. Investigaremos y seleccionaremos una variedad de instrumentos musicales virtuales que se ajusten a las temáticas predefinidas de los videojuegos. Configuraremos estos instrumentos para que reproduzcan las melodías generadas de manera convincente al oído humano.
    \item Desarrollo de la aplicación de escritorio. Utilizaremos las tecnologías adecuadas para desarrollar una interfaz de usuario intuitiva y amigable que permita a los usuarios generar, armonizar y sonorizar sus propias composiciones musicales para videojuegos.
\end{enumerate}








\chapter*{Conclusions and Future Work}
\label{cap:conclusions}
\addcontentsline{toc}{chapter}{Conclusions and Future Work}

Conclusions and future lines of work. This chapter contains the translation of Chapter \ref{cap:conclusiones}.



\end{otherlanguage}
%%%%%%%%%%%%%%%%%%%%%%%%%%%%%%%%%%%%%%%%%%%%%%%%%%%%%%%%%%%%%%%%%%%%%%%%%%%

\chapter*{Contribuciones Personales}
\label{cap:contribucionesPersonales}
\addcontentsline{toc}{chapter}{Contribuciones Personales}

\section*{Rodrigo Sánchez Torres}
\subsection*{Procesamiento de datos}
Para poder utilizar algoritmos de inteligencia artificial en nuestro trabajo era necesario tener un dataset limpio y adecuado para nuestros objetivos.

La primera tarea fue buscar un dataset que contenga información sobre melodías. En la búsqueda se valoraron muchos tipos de datasets, finalmente, en la exploración de los datasets de Magenta se decidió utilizar el \textit{Bach Doodle Dataset}.

Una vez elegido el dataset se procedió al tratamiento de este, con distintos acercamientos se fueron realizando limpiezas en principio superficiales, pero pronto se decidió que lo mejor era tener un dataset con solamente melodías que estén en la misma escala e incluso en un rango reducido de octavas. 

También se exploraron las opciones de usar o no silencios, siendo en algoritmos como cadenas de Markov mejor no emplear silencios, pero en redes neuronales recurrentes se obtuvieron mejores resultados con silencios en el dataset.

Finalmente se testeó tanto a usar tiempos absolutos en segundos como a normalizar el tiempo utilizando \textit{steps} y BPM (\textit{beats per minute}), dando mejor resultado tener el tiempo normalizado y homogéneo en el dataset.

\subsection*{Algoritmos de \textit{Machine Learning}}
Antes de poder elegir qué modelos serían más adecuados para la generación de melodías fue necesario realizar una amplia investigación sobre el tema, para poder conocer la mayor cantidad de técnicas y modelos que podemos aplicar a nuestro trabajo.

Tras la investigación quedó claro que era necesario buscar un modelo que no fuera determinista, para poder generar melodías diferentes en cada ejecución. Es por esto que se decidió explorar los 3 modelos actuales: cadenas de Markov, redes neuronales recurrentes con capa \textit{softmax} y modelos de Magenta.

En primer lugar se empezó a explorar la generación con Markov, para esto se realizó una investigación sobre el tema y sus posibles aplicaciones. Para poder utilizar correctamente las cadenas se utilizó una notación común a lo largo del proyecto en los algoritmos propios de generación, como era \textit{"pitch\_duration"}.

Para poder utilizar cadenas de Markov se procesó manualmente el dataset, creando matrices de transición que sean interpretables por la librería utilizada. Finalmente, se consiguió tener un primer modelo de generación listo, que permitió continuar el desarrollo de otras ramas que se apoyan en éste.

Una vez se tuvo el primer modelo de generación listo se podía dedicar más tiempo a desarrollar un modelo más complejo. Se decidió utilizar entonces una red neuronal recurrente. Pero esta tenía el problema de ser un modelo determinista, así que se añadió finalmente una capa \textit{softmax}, que permitía introducir una aleatoriedad a la generación controlada con un parámetro de temperatura.

Este modelo fue el más complejo de desarrollar, pues había que realizar mucho trabajo de ajuste de \textit{súperparámetros} y diseño de la propia red. Además, se realizaron varias modificaciones al dataset, que finalmente no aportaban un mejor resultado. Al final tras un largo período de prueba ser decidió adoptar el diseño actual, teniendo una entrada de número reales, con el \textit{pitch}, duración y tiempo de inicio de la nota actual, y una salida con valores de tipo \textit{label}, siendo estas las posibles notas que podía elegir la red. Al utilizar salida con valores de tipo \textit{label} codificados mediantes \textit{one-hot encoding} se podía utilizar la capa \textit{softmax} mencionada anteriormente.

La red conseguía finalmente un rendimiento aceptable, por lo que se realizaron pruebas para entrenarla de forma que indicara también el tiempo de inicio de cada nota, lo que podía generar melodías más interesantes sobre el papel. Sin embargo, la red no conseguía aprender patrones significativos sobre este, obteniendo una precisión tan baja que no merecía la pena utilizar dicho parámetro. Es por esto que se decidió introducir los silencios como notas adicionales. 

Este modelo consiguió una generación mucho mejor a las cadenas de Markov, manteniendo patrones más típicos en la composición pero con limitaciones. Aún así se decidió mantener el modelo con Markov para tener una generación más experimental.

Intercalado con el desarrollo de la red neuronal recurrente se integró Magenta al proyecto, con el fin de proporcionar una generación más elaborada que la de Markov para poder trabajar en otras ramas y producir mejores resultados. 

La integración de Magenta supuso una gran cantidad de problemas, pues principalmente no se podía utilizar el paquete de Python en Windows. Tras realizar varias pruebas con binarios precompilados para Windows, migraciones a Ubuntu y diversos acercamientos se decidió que no era viable utilizar el paquete de Python, por lo que se exploró el uso del paquete de JavaScript.

El principal conflicto era que el proyecto estaba íntegramente desarrollado en Python, por lo que el módulo de JavaScript tenía que estar separado de éste. Finalmente se consiguió comunicar el proyecto con un script (magentaPython.py) que actúa de puente entre Python y JavaScript. Este script utiliza el módulo \textit{subproccess} para lanzar un proceso de JavaScript, con el que nos podemos comunicar por medio de argumentos de dicho proceso y la salida estándar de éste. Con esta metodología podemos mandar melodías serializadas en JSON al proceso y recibir resultados de generación recogiendo la salida estándar con el mismo formato JSON.

Este acercamiento permite utilizar Magenta en la aplicación, pero conlleva otros problemas de dependencia de paquetes y sobre todo de una conexión a internet para utilizar los modelos alojados en la base de datos de Google.

\subsection*{Aportaciones a la aplicación}
Además, se han realizado aportaciones al código de la aplicación. Se han introducido los módulos de generación a ésta, así como la posibilidad de elegir uno u otro desde código. También se han añadido las funcionalidades de carga de melodías existentes y guardado de una melodía generada, pudiendo elegir el nombre y ubicación del archivo.

\section*{Víctor Manuel Estremera Herranz}

Antes de empezar con las técnicas de armonización utilizadas, falta un paso previo crucial. Y es la creación de una estructura de clases y métodos que den soporte a todo lo explicado en el apartado \ref{sec:arm:armonia}. Existen clases que abstraen lo que representa una nota musical (\ref{arm:notas_musicales}), un intervalo (\ref{sec:arm:intervalos}), una escala (\ref{sec:arm:escalas}) y la armonía de una escala (\ref{arm:armonia_escala}) tal como se ha explicado. Un acorde (\ref{sec:arm:acordes}) es respresentado también como una escala. 

Al menos dos páginas con las contribuciones del estudiante 2. En caso de que haya más estudiantes, copia y pega una de estas secciones.

\section*{Javier Callejo Herrero}
\subsection*{Generación de percusiones.}
Lo primero que se hizo fue investigar los diferentes modelos de Magenta, tanto para tomar ideas como para decidir si podíamos usar dichos modelos en nuestra herramienta. 

Una vez vistas las herramientas disponibles, se decidió que los resultados obtenidos no ofrecían la variedad que requería la herramienta, por lo que se implementó un generador de percusión propio. Este generador permite crear patrones de batería de varios géneros distintos (como el jazz, el rock, la música disco, etc.) Estos patrones generados nos servirán para las siguientes partes del trabajo, en concreto la del diseño de temáticas.

\subsection*{Generación de líneas de bajo.}
Tras haber implementado el generador de percusión, se implementó una primera versión de un generador de líneas de bajo, que cogía un ítem MIDI con acordes y usaba las notas más graves para crear de forma pseudoaleatoria una melodía monofónica que si  sirviera como línea de bajo. Más adelante se descartó el uso de este generador al descubrir formas mejores de afrontar este programa (el uso de un plugin arpegiador).

\subsection*{Selección de plugins.}
De forma paralela a los apartados anteriores, se realiza una investigación exhaustiva de plugins (de instrumentos virtuales principalmente). Continuamente nos encontramos con problemas de almacenamiento e instalación muy tediosa.

Es más adelante cuando se reúnen finalmente todos los plugins de instrumentos virtuales que se usarán en la herramienta, plugins ligeros, de fácil instalación, pues son colocados los .vst o .dll en una única carpeta comprimida para facilitar la instalación tanto al usuario final como a los propios compañeros del trabajo.

\subsection*{Programar en Reaper.}
Se investiga el uso de ReaScript con Python para comenzar a automatizar tareas en Reaper. Primeramente se logran cargar instrumentos virtuales e ítems MIDI llamando a un script desde una acción de Reaper.

A continuación, se diseña y se implementa la carga de plugins que cubren las necesidades concretas de cada pista. (Se continúa buscando plugins que cubran nuestras necesidades, en particular plugins de efectos).

Desde ahí, se comienza a desarrollar una arquitectura que permite crear un arreglo para la primera temática (pradera). En este punto tenemos un script que carga ítems MIDI de melodía, armonía, batería y bajo. Además, añade de forma pseudoaleatoria plugins a las pistas y al máster: tenemos las primeras canciones generadas.

En este momento se empiezan a diseñar las temáticas. Se habla de ello en el siguiente apartado.

Se continúa hasta tener todas las temáticas programadas, añadiendo progresivamente mejoras y funcionalidades nuevas, como el subir o bajar semitonos, añadir diversos efectos a la mezcla o, finalmente, cambiar la complejidad de la melodía cargada.

Cabe destacar que existe una investigación sobre conectar scripts externos con Reaper, pero no se logra dicho objetivo y pasa a un segundo plano hasta más avanzado el desarrollo (Sección de Miguel Comunicación entre Vanguard Music y Reaper).

\subsection*{Diseño e implementación de las temáticas.}
Se investigan los elementos tanto técnicos como creativos necesarios para crear las distintas temáticas para las canciones. A continuación, se diseñan e implementan en Reaper todas las temáticas que están disponibles en la herramienta.

Para la implementación de las temáticas, ha sido requerido un exhaustivo trabajo de selección de timbres usando los instrumentos virtuales disponibles. El proceso es el siguiente: selección de plugin, elección o diseño del timbre dentro del plugin, guardado de un preset de Reaper y adición de dicho preset al script de Python para su posterior recuperación al lanzar el script en Reaper.

\subsection*{Diseño de la interfaz gráfica.}
Finalmente, se renderizaron en 3D usando Blender los diversos fondos presentes en la pestaña de sonorización de la aplicación y se colocaron en TKInter.


\section*{Miguel González Pérez}

\subsection*{Primera aproximación a la generación de MIDI con Magenta.}
Uno de los primeros objetivos era conseguir generar melodías con música simbólica a través de la herramienta de Magenta. Dado que no fue posible para nosotros utilizar la API de Magenta desde Python en Windows, surgió la alternativa de utilizar la API de JavaScript, lo que implicaba integrar Node.js con el resto del desarrollo en Python. 

Si bien logramos adquirir la melodía en forma de NoteSequence de Magenta en el contexto de JavaScript, no podíamos acceder a ello desde Python, por lo que más tarde se encargaría mi compañero Rodrigo de continuar y completar esta tarea.

\subsection*{Arquitectura de la Aplicación en TKInter.}
Tras el anterior apartado, mi siguiente tarea fue la de montar la arquitectura de la aplicación en TKInter con Python. Gran parte del trabajo en este apartado fue el del estudio de la documentación de \textit{TKInter}, la biblioteca estándar de Python para la creación de interfaces gráficas. 

Utilizando estos conocimientos, diseñé la aplicación de forma que disgregase las funcionalidades principales en pestañas diferenciadas. Esto requería en cuanto a implementación diseñar el código de forma que gestionar y añadir varias pestañas fuera escalable y sostenible.

El mantenimiento de un estado interno de la aplicación durante la ejecución y la posterior serialización a JSON y persistencia en ficheros en forma de presets son parte del trabajo que he realizado en el desarrollo de la aplicación. De este estado interno se sustrae la información necesaria, como temáticas y semillas, para la instrumentalización en Reaper. Esto implica también recuperar las configuraciones del estado interno de la aplicación guardadas por el usuario, de forma que no produzca errores.

La configuración propia de la aplicación, como la ruta al ejecutable de Reaper, también fue parte de este proceso de persistencia en formato JSON.

Para el desarrollo de la aplicación he aplicado algunos patrones de diseño software, como \textit{Strategy} a la hora de asignar el generador musical en la pestaña \textit{avanzado} o el propio sistema de pestañas de la aplicación, que aplica el patrón \textit{State} para manejar la entrado, actualización y salida de cada módulo de la aplicación.

En un momento avanzado del desarrollo de la aplicación, encontramos la necesidad de explicar los distintos elementos de la interfaz de usuario de la aplicación de forma rápida y sencilla. Para esto se creó la clase \textit{Tooltip} que permitía mostrar un pequeño texto cuando pasabas el ratón por encima de un elemento de la interfaz. Más tarde, cuando la aplicación evolucionó, optamos por usar un plugin de TKInter que traía \textit{tooltips} integradas que funcionaban mejor, adaptándose al tema oscuro de la aplicación.

\subsection*{Comunicación entre Vanguard Music y Reaper.}
Uno de los principales retos de la aplicación era la comunicación con Reaper. Nuestra intención inicial era la de permitir al usuario utilizar la aplicación sin tener que preocuparse en exceso por Reaper. Esto planteaba varios desafíos, siendo uno de ellos el de añadir nuestros scripts de ejecución sin necesidad de que el usuario tenga que registrarlos manualmente en Reaper. 

Reaper permite la ejecución de scripts en Python, \textit{LUA} y \textit{EEL} a través de su API \textit{ReaScript}. No obstante, estos scripts tienen que haber sido registrados en Reaper antes de poder ser ejecutad. Si bien éramos capaces de ejecutar los scripts ya registrados en Reaper desde nuestra aplicación gracias al trabajo de Javi, no podíamos registrar nuevos scripts desde código. Mi trabajo en esta sección fue la de estudiar la documentación de \textit{ReaScript} con el fin de enviar nuestros scripts a Reaper desde código. 

Los scripts registrados en Reaper tienen un ID único asociado, no obstante este identificador hemos visto que se puede perder en algunos casos, por ejemplo, en cambios de versiones de Reaper, por lo que hacemos ahora es registrar de nuevo los scripts cuando lanzamos Reaper desde la aplicación, sustituyendo los scripts homónimos que hubiese registrado antes en el proceso.






%
% Bibliografía
%
% Si el TFM se escribe en inglés, editar TeXiS/TeXiS_bib para cambiar el
% estilo de las referencias
%---------------------------------------------------------------------
%
%                      configBibliografia.tex
%
%---------------------------------------------------------------------
%
% bibliografia.tex
% Copyright 2009 Marco Antonio Gomez-Martin, Pedro Pablo Gomez-Martin
%
% This file belongs to the TeXiS manual, a LaTeX template for writting
% Thesis and other documents. The complete last TeXiS package can
% be obtained from http://gaia.fdi.ucm.es/projects/texis/
%
% Although the TeXiS template itself is distributed under the 
% conditions of the LaTeX Project Public License
% (http://www.latex-project.org/lppl.txt), the manual content
% uses the CC-BY-SA license that stays that you are free:
%
%    - to share & to copy, distribute and transmit the work
%    - to remix and to adapt the work
%
% under the following conditions:
%
%    - Attribution: you must attribute the work in the manner
%      specified by the author or licensor (but not in any way that
%      suggests that they endorse you or your use of the work).
%    - Share Alike: if you alter, transform, or build upon this
%      work, you may distribute the resulting work only under the
%      same, similar or a compatible license.
%
% The complete license is available in
% http://creativecommons.org/licenses/by-sa/3.0/legalcode
%
%---------------------------------------------------------------------
%
% Fichero  que  configura  los  parámetros  de  la  generación  de  la
% bibliografía.  Existen dos  parámetros configurables:  los ficheros
% .bib que se utilizan y la frase célebre que aparece justo antes de la
% primera referencia.
%
%---------------------------------------------------------------------


%%%%%%%%%%%%%%%%%%%%%%%%%%%%%%%%%%%%%%%%%%%%%%%%%%%%%%%%%%%%%%%%%%%%%%
% Definición de los ficheros .bib utilizados:
% \setBibFiles{<lista ficheros sin extension, separados por comas>}
% Nota:
% Es IMPORTANTE que los ficheros estén en la misma línea que
% el comando \setBibFiles. Si se desea utilizar varias líneas,
% terminarlas con una apertura de comentario.
%%%%%%%%%%%%%%%%%%%%%%%%%%%%%%%%%%%%%%%%%%%%%%%%%%%%%%%%%%%%%%%%%%%%%%
\setBibFiles{%
biblio%
}

%%%%%%%%%%%%%%%%%%%%%%%%%%%%%%%%%%%%%%%%%%%%%%%%%%%%%%%%%%%%%%%%%%%%%%
% Definición de la frase célebre para el capítulo de la
% bibliografía. Dentro normalmente se querrá hacer uso del entorno
% \begin{FraseCelebre}, que contendrá a su vez otros dos entornos,
% un \begin{Frase} y un \begin{Fuente}.
%
% Nota:
% Si no se quiere cita, se puede eliminar su definición (en la
% macro setCitaBibliografia{} ).
%%%%%%%%%%%%%%%%%%%%%%%%%%%%%%%%%%%%%%%%%%%%%%%%%%%%%%%%%%%%%%%%%%%%%%
\setCitaBibliografia{
\begin{FraseCelebre}
\begin{Frase}
  Y así, del mucho leer y del poco dormir, se le secó el celebro de
  manera que vino a perder el juicio.\\ 
  \textcolor{red}{(modificar en Cascaras$\backslash$bibliografia.tex)}
\end{Frase}
\begin{Fuente}
  Miguel de Cervantes Saavedra
\end{Fuente}
\end{FraseCelebre}
}

%%
%% Creamos la bibliografia
%%
\makeBib

% Variable local para emacs, para  que encuentre el fichero maestro de
% compilación y funcionen mejor algunas teclas rápidas de AucTeX

%%%
%%% Local Variables:
%%% mode: latex
%%% TeX-master: "../Tesis.tex"
%%% End:



% Apéndices
\appendix
\chapter{Dependencias de Instalación}
\label{Appendix:Key1}

Para poder hacer uso nuestra aplicación, se requiere tener instaladas herramientas de terceros. Estas herramientas son esenciales para el funcionamiento correcto de nuestra aplicación.
Cada dependencia tiene asociado el enlace donde se puede descargar este contenido, véase


\begin{table}
    \centering
	\begin{tabular}{c|c|c|c}
		\textbf{Herramienta} & \textbf{Versión recomendada} & \textbf{Enlace} & \textbf{Fecha Enlace} \\
		\hline\hline
		python & 3.9 & \href{https://www.python.org/downloads/release/python-390/}{downloads} & \today\\
		6 & 2.12 & 4.40 & \today\\
		1 & 3.79 & 5.00 & \today\\
		2 & 4.88 & 5.30 & \today\\
		4 & 3.50 & 2.90 & \today\\
		5 & 7.40 & 4.70 & \today\\
		\hline
	\end{tabular}
	\caption{Tabla de ejemplo}
	\label{tab:sampleTable}
\end{table}
\chapter{Representaciones de una canción}
\label{Appendix:Key2}

    En el módulo al que se hace referencia durante toda la Sección \ref{cap:armonizacion} se manejan dos representaciones distintas de las notas:

    \begin{enumerate}
        \item \textbf{Lista de notas}: es la representación estándar en este módulo de la aplicación. Es la más cómoda para realizar operaciones simples con las notas, como por ejemplo recortar una canción, transponer todas las notas, etc... Se almacenan en una lista desordenada de estructuras con el siguiente formato:
        \begin{enumerate}
            \item[\textbullet] \textbf{note}: almacena el \textit{pitch} MIDI (tono) de la nota. Guardar el \textit{pitch} MIDI es similar a guardar la nota musical, ya que este almacena de forma implícita el nombre de la nota y su octava.
                       
\begin{table}[htbp]
    \centering
    \begin{tabular}{c||c|c|c|c|c|c|c|c|c|c|c|c}
        \textbf{Nombre} & C & Db & D & Eb & E & F & Gb & G & Ab & A & Bb & B \\ 
        \hline
        \textbf{Pitch} &  0 & 1 & 2 & 3 & 4 & 5 & 6 & 7 & 8 & 9 & 10 & 11 \\
    \end{tabular}
    \caption{Relación entre nota y su \textit{pitch} en la primera octava}
\label{tab:nota_pitch}
\end{table}

         Por ejemplo, si calculamos a qué nota le corresponde el \textit{pitch} 40, sabemos que el nombre de la nota es E si nos fijamos en la tabla, ya que 40 mod 12 = 4 y sabemos que pertenece a cuarta octava, ya que 40 div 12 = 3 (se empieza en la octava 0).

        \item[\textbullet] \textbf{start\_time}: tiempo en \textit{ticks} en el que empieza a sonar una nota desde que empieza la canción en el \textit{tick} 0.
        \item[\textbullet] \textbf{duration}: tiempo en \textit{ticks} desde que empieza a sonar la nota hasta que para.
    \end{enumerate}
    \item \textbf{Diccionario de eventos}: se almacena en cada \textit{tick} clave el evento que ha ocurrido. Los eventos están ordenados de menor a mayor según el \textit{tick} en el que ocurren. Como tal, solo existen dos tipos de eventos: \textbf{note\_on} y \textbf{note\_off}.

    A cada evento viene asociado el \textit{pitch} de la nota afectada. Esto puede suponer, a priori, un inconveniente, ya que, si hay dos notas del mismo \textit{pitch} superpuestas en el mismo espacio de tiempo, existe una ambigüedad a la hora de saber qué evento \textit{note\_off} le corresponde a cada una. Sin embargo, como esta representación es únicamente utilizada en la armonización por ventanas (Sección \ref{arm:sec:ventanas_normal}) para hacer un recorrido lineal de las notas de la melodía, este inconveniente no les afecta. 

    Existen métodos para pasar de la anterior representación a esta, pero no al revés. De nuevo, esto se debe a que esta representación solo se utiliza para realizar la armonización por ventanas.
\end{enumerate}
%\include{Apendices/appendixC}
%\include{...}
%\include{...}
%\include{...}
\backmatter



%
% Índice de palabras
%

% Sólo  la   generamos  si  está   declarada  \generaindice.  Consulta
% TeXiS.sty para más información.

% En realidad, el soporte para la generación de índices de palabras
% en TeXiS no está documentada en el manual, porque no ha sido usada
% "en producción". Por tanto, el fichero que genera el índice
% *no* se incluye aquí (está comentado). Consulta la documentación
% en TeXiS_pream.tex para más información.
\ifx\generaindice\undefined
\else
%%---------------------------------------------------------------------
%
%                        TeXiS_indice.tex
%
%---------------------------------------------------------------------
%
% TeXiS_indice.tex
% Copyright 2009 Marco Antonio Gomez-Martin, Pedro Pablo Gomez-Martin
%
% This file belongs to TeXiS, a LaTeX template for writting
% Thesis and other documents. The complete last TeXiS package can
% be obtained from http://gaia.fdi.ucm.es/projects/texis/
%
% This work may be distributed and/or modified under the
% conditions of the LaTeX Project Public License, either version 1.3
% of this license or (at your option) any later version.
% The latest version of this license is in
%   http://www.latex-project.org/lppl.txt
% and version 1.3 or later is part of all distributions of LaTeX
% version 2005/12/01 or later.
%
% This work has the LPPL maintenance status `maintained'.
% 
% The Current Maintainers of this work are Marco Antonio Gomez-Martin
% and Pedro Pablo Gomez-Martin
%
%---------------------------------------------------------------------
%
% Contiene  los  comandos  para  generar  el índice  de  palabras  del
% documento.
%
%---------------------------------------------------------------------
%
% NOTA IMPORTANTE: el  soporte en TeXiS para el  índice de palabras es
% embrionario, y  de hecho  ni siquiera se  describe en el  manual. Se
% proporciona  una infraestructura  básica (sin  terminar)  para ello,
% pero  no ha  sido usada  "en producción".  De hecho,  a pesar  de la
% existencia de  este fichero, *no* se incluye  en Tesis.tex. Consulta
% la documentación en TeXiS_pream.tex para más información.
%
%---------------------------------------------------------------------


% Si se  va a generar  la tabla de  contenidos (el índice  habitual) y
% también vamos a  generar el índice de palabras  (ambas decisiones se
% toman en  función de  la definición  o no de  un par  de constantes,
% puedes consultar modo.tex para más información), entonces metemos en
% la tabla de contenidos una  entrada para marcar la página donde está
% el índice de palabras.

\ifx\generatoc\undefined
\else
   \addcontentsline{toc}{chapter}{\indexname}
\fi


% Generamos el índice
\printindex

% Variable local para emacs, para  que encuentre el fichero maestro de
% compilación y funcionen mejor algunas teclas rápidas de AucTeX

%%%
%%% Local Variables:
%%% mode: latex
%%% TeX-master: "./tesis.tex"
%%% End:

\fi

%
% Lista de acrónimos
%

% Sólo  lo  generamos  si  está declarada  \generaacronimos.  Consulta
% TeXiS.sty para más información.


\ifx\generaacronimos\undefined
\else
%---------------------------------------------------------------------
%
%                        TeXiS_acron.tex
%
%---------------------------------------------------------------------
%
% TeXiS_acron.tex
% Copyright 2009 Marco Antonio Gomez-Martin, Pedro Pablo Gomez-Martin
%
% This file belongs to TeXiS, a LaTeX template for writting
% Thesis and other documents. The complete last TeXiS package can
% be obtained from http://gaia.fdi.ucm.es/projects/texis/
%
% This work may be distributed and/or modified under the
% conditions of the LaTeX Project Public License, either version 1.3
% of this license or (at your option) any later version.
% The latest version of this license is in
%   http://www.latex-project.org/lppl.txt
% and version 1.3 or later is part of all distributions of LaTeX
% version 2005/12/01 or later.
%
% This work has the LPPL maintenance status `maintained'.
% 
% The Current Maintainers of this work are Marco Antonio Gomez-Martin
% and Pedro Pablo Gomez-Martin
%
%---------------------------------------------------------------------
%
% Contiene  los  comandos  para  generar  el listado de acrónimos
% documento.
%
%---------------------------------------------------------------------
%
% NOTA IMPORTANTE:  para que la  generación de acrónimos  funcione, al
% menos  debe  existir  un  acrónimo   en  el  documento.  Si  no,  la
% compilación  del   fichero  LaTeX  falla  con   un  error  "extraño"
% (indicando  que  quizá  falte  un \item).   Consulta  el  comentario
% referente al paquete glosstex en TeXiS_pream.tex.
%
%---------------------------------------------------------------------


% Redefinimos a español  el título de la lista  de acrónimos (Babel no
% lo hace por nosotros esta vez)

\def\listacronymname{Lista de acrónimos}

% Para el glosario:
% \def\glosarryname{Glosario}

% Si se  va a generar  la tabla de  contenidos (el índice  habitual) y
% también vamos a  generar la lista de acrónimos  (ambas decisiones se
% toman en  función de  la definición  o no de  un par  de constantes,
% puedes consultar config.tex  para más información), entonces metemos
% en la  tabla de contenidos una  entrada para marcar  la página donde
% está el índice de palabras.

\ifx\generatoc\undefined
\else
   \addcontentsline{toc}{chapter}{\listacronymname}
\fi


% Generamos la lista de acrónimos (en realidad el índice asociado a la
% lista "acr" de GlossTeX)

\printglosstex(acr)

% Variable local para emacs, para  que encuentre el fichero maestro de
% compilación y funcionen mejor algunas teclas rápidas de AucTeX

%%%
%%% Local Variables:
%%% mode: latex
%%% TeX-master: "../Tesis.tex"
%%% End:

\fi

%
% Final
%
%%---------------------------------------------------------------------
%
%                      fin.tex
%
%---------------------------------------------------------------------
%
% fin.tex
% Copyright 2009 Marco Antonio Gomez-Martin, Pedro Pablo Gomez-Martin
%
% This file belongs to the TeXiS manual, a LaTeX template for writting
% Thesis and other documents. The complete last TeXiS package can
% be obtained from http://gaia.fdi.ucm.es/projects/texis/
%
% Although the TeXiS template itself is distributed under the 
% conditions of the LaTeX Project Public License
% (http://www.latex-project.org/lppl.txt), the manual content
% uses the CC-BY-SA license that stays that you are free:
%
%    - to share & to copy, distribute and transmit the work
%    - to remix and to adapt the work
%
% under the following conditions:
%
%    - Attribution: you must attribute the work in the manner
%      specified by the author or licensor (but not in any way that
%      suggests that they endorse you or your use of the work).
%    - Share Alike: if you alter, transform, or build upon this
%      work, you may distribute the resulting work only under the
%      same, similar or a compatible license.
%
% The complete license is available in
% http://creativecommons.org/licenses/by-sa/3.0/legalcode
%
%---------------------------------------------------------------------
%
% Contiene la última página
%
%---------------------------------------------------------------------


% Ponemos el marcador en el PDF
\ifpdf
   \pdfbookmark{Fin}{fin}
\fi

\thispagestyle{empty}\mbox{}

Este texto se puede encontrar en el fichero Cascaras/fin.tex. Si deseas eliminarlo, basta con comentar la línea correspondiente al final del fichero TFGTeXiS.tex.

\vspace*{4cm}

\small

\hfill \emph{--¿Qué te parece desto, Sancho? -- Dijo Don Quijote --}

\hfill \emph{Bien podrán los encantadores quitarme la ventura,}

\hfill \emph{pero el esfuerzo y el ánimo, será imposible.}

\hfill 

\hfill \emph{Segunda parte del Ingenioso Caballero} 

\hfill \emph{Don Quijote de la Mancha}

\hfill \emph{Miguel de Cervantes}

\vfill%space*{4cm}

\hfill \emph{--Buena está -- dijo Sancho --; fírmela vuestra merced.}

\hfill \emph{--No es menester firmarla -- dijo Don Quijote--,}

\hfill \emph{sino solamente poner mi rúbrica.}

\hfill 

\hfill \emph{Primera parte del Ingenioso Caballero} 

\hfill \emph{Don Quijote de la Mancha}

\hfill \emph{Miguel de Cervantes}


\newpage
\thispagestyle{empty}\mbox{}

\newpage

% Variable local para emacs, para  que encuentre el fichero maestro de
% compilación y funcionen mejor algunas teclas rápidas de AucTeX

%%%
%%% Local Variables:
%%% mode: latex
%%% TeX-master: "../Tesis.tex"
%%% End:

%\end{otherlanguage}
\end{document}
